% ---
% id: meat-core/rem:declared-measurement
% kind: rem
% title: "The declared measurement as a premise of the rate"
% summary: >-
%   A tiny change of the declared measurement can move the key rate by many
%   bits per round, so the model is a premise of the rate. Where a handed
%   reference's companion counts and the key length is formed on it, the snap
%   only steers the search; nothing bounds model to device.
% user_statement: >-
%   The key rate is computed from the measurement model you declare for your
%   receiver, and that model is a premise of the number: nothing checks that
%   your device measures what it says (for a receiver the engine models,
%   [[A-TRUST]]). A change of the model too small to see in its entries,
%   notably one that makes an outcome nearly or exactly impossible, can move
%   the rate by many bits per round in either direction; no result bounds
%   that gap. The engine also moves your model onto its exact number grid (the
%   snap) for its search. Where your model reaches the engine as a reference,
%   the bound is evaluated on your model itself and needs no premise about the
%   snap, unless the rate is certified only on the moved model, as the report
%   then says; any other block rests on [[A-SNAP]], or the run refuses if it
%   declares the strict option. Validate the model, above all which outcomes
%   it treats as impossible.
% prominence: standard
% status: accepted
% uses: [prop:optol-nobound, prop:value-range, def:declaration, rem:optol-direction, req:infinite-value, prop:reference-companion, foundations/def:snap, foundations/req:snap-reference, foundations/rem:noise-floor, foundations/def:evaluated-off-snap]
% mentions: [thm:keylength, foundations/ass:snapped-measurement, receiver/ass:trust, meat-core/frank-wolfe/cor:fo-reference, meat-core/frank-wolfe/req:bracket-on-reference, meat-core/frank-wolfe/req:certified-evaluation, meat-core/frank-wolfe/rem:fo-arithmetic, req:companion-admission, reductions/req:reference-companion]
% assumes: []
% refs: []
% provenance: newly-derived
% lean:
% ---

A key rate of this contract is computed from a declaration whose
measurement on $\Qsys$ is a model of a device, and
Proposition~\ref{prop:optol-nobound} shows how much the number rests on that
model. A domination weight declared for the device, outright or derived from
ranges of its efficiencies and dark-count probabilities, would carry part of
the gap between model and device into the flags of a \tref{domination link}; it is
future scope (\repodoc{spec/future-scope.md}~\S13).

\begin{remark}\label{rem:declared-measurement}
The value $\hlow$ is not controlled by a tolerance on the declared
measurement $\{F_v\}$ on $\Qsys$ (Definition~\ref{def:declaration}) at any
tolerance down to $2^{-52}$. For every $\epsilon\in(0,\tfrac12]$ and
every positive integer $N$ with $2^{N}\le\epsilon\,2^{53}$ there is a
declaration on which two \tref{exactly admissible} families within
$\epsilon$ of each other in every element give values
$\frac{\alpha}{\alpha-1}N$ bits per round apart, and no allowance computed
from the declared dimensions, alphabet sizes and $\alpha$ below that scale
bounds the change (Proposition~\ref{prop:optol-nobound}\,(i)--(ii)); below
$\epsilon=2^{-52}$ the proposition exhibits no such pair. For
$\alpha\in(1,2)$ the gap reaches the key length of
Theorem~\ref{thm:keylength}: over $n$ rounds the two lengths differ by at
least $nN-1$ bits whenever the smaller is positive
(Proposition~\ref{prop:optol-nobound}\,(iii)). The gap is bought by a ratio,
not by a size: on the instance of that proposition
$\hlow=\frac{\alpha}{\alpha-1}\log_2\frac1{g_0}$, where $g_0$ is the
probability of the one outcome the acceptance test demands, and moving
$g_0$ from $\epsilon'$, the largest power of two at most $\epsilon$, to
$\epsilon'2^{-N}$ changes the element by at most $\epsilon$ and the
probability by the factor $2^{N}$. Its limit is a change of support: at
$g_0=0$ the measurement treats the demanded outcome as impossible, no
feasible state gives it positive probability, and $\hlow=+\infty$
(Proposition~\ref{prop:value-range}\,(iv)), a value that carries no
key-length claim (Requirement~\ref{req:infinite-value}). A change of which
outcomes the measurement treats as impossible, or nearly so, can thus move
the rate by far more than the size of the change, in either direction; it
need not, since the proposition exhibits one declaration on which it does
and states no property of every declaration
(Remark~\ref{rem:optol-direction}). The declared measurement model is
therefore itself a premise of every number computed from it.

Between the device and the declared measurement lie two moves. The first is
the \tref{snap} of a family read as a measurement onto its lattice, whose
displacement, however small, bounds the change of the rate by nothing.
Where the declarer holds a \tref{reference} of the family and hands it
beside the snap, a claim establishes for the instance that the construction
may run on it (\xref{reductions/req:reference-companion}), the companion is
compiled and admitted (Requirement~\ref{req:companion-admission}), and the
run's bound on the companion counts
(Requirement~\ref{req:bracket-on-reference}\,(iii)) and its key length is
formed from a number on the companion, not from one certified on the snapped
declaration in the fallback of
Requirement~\ref{req:certified-evaluation}\,(i), the run need
not bound that change at all. Its first-order bound is then evaluated on the compiled reference companion
(Proposition~\ref{prop:reference-companion}), at the points the search on the
snapped declaration produced and at those of a short search continued on the
companion from them (Corollary~\ref{cor:fo-reference},
Requirement~\ref{req:bracket-on-reference}); the snap only steers the choice
of those points, and a bound that holds at any admissible point holds at
them. Every
block on which the companion carries the reference is
\tref{evaluated off its snap} (\xref{foundations/def:evaluated-off-snap}),
nothing the key length rests on reads its snapped data, and it rests on no
premise about the snap. In that fallback the number the key length is formed
from reads the snapped data of those blocks, which are then among the blocks
the run does not evaluate off its snap, below
(Requirement~\ref{req:bracket-on-reference}\,(v)). Outside that fallback,
the key length is one of the model the reference describes, and what lies between it and
the device is the second move below. The evaluation of an enclosed
reference at its centres is itself a displacement of the measured family, of
the size of the enclosure's radius, to which
Proposition~\ref{prop:optol-nobound} applies as it does to the snap's: the
premise about the snap is gone, the exposure to the radius is not, and the
recorded largest radius is the only handle on it
(Remark~\ref{rem:fo-arithmetic}). The rest of this paragraph concerns the
blocks a run does not evaluate off its snap. A block kept because its
reference is exact and already on the lattice rests on nothing, the snap
being the identity there. Every other block is a \tref{moved} block of the
snap (\xref{foundations/def:snap}), and no claim relates it to the
reference: domination, which would bound the snapped family by the
reference, is not established for a snap, and an idealised model may
dominate no lattice family at a small weight at all
(\xref{foundations/rem:noise-floor}). On such a block the run proceeds,
unless even the last resort yields no snap, resting on
\xref{foundations/ass:snapped-measurement}: that on the blocks its record
names the device measures the snapped family, or, where the snapped family is
a producer's emission, that what the run concludes from the emission holds
of the snapped family. No claim then relates its key length to the
reference, and a run that declares the strict option refuses instead
(\xref{foundations/req:snap-reference}\,(iv)). The second move is
the gap between the model as declared, or its reference, and the physical
device. No claim bounds it: the run assumes, without checking, that the
model describes the device, for a receiver that the receiver module models
through \xref{receiver/ass:trust}, and by
Proposition~\ref{prop:optol-nobound} a small distance between model and
device bounds the change of the rate no more than a small distance between
two declared families does.
\end{remark}
